% Initial research-talk deck based on overleaf/main_linear.tex.
% Remove `handout` to enable overlays when presenting.
\documentclass[handout,compress,9pt]{beamer}

\usefonttheme[onlymath]{serif}
\usepackage[english]{babel}
\usepackage{amsmath,amssymb,amsthm}
\usepackage{bm}
\usepackage{booktabs}
\usepackage{graphicx}
\usepackage{tikz}
\usetikzlibrary{arrows.meta,positioning,calc}

\usetheme{Boadilla}
\useoutertheme{infolines}
\setbeamertemplate{navigation symbols}{}
\setbeamertemplate{theorems}[numbered]
\setbeamertemplate{caption}[numbered]

\definecolor{ddlblue}{RGB}{0,78,140}
\definecolor{ddlred}{RGB}{176,34,34}
\definecolor{ddlgreen}{RGB}{32,120,82}
\newcommand{\blue}[1]{\textcolor{ddlblue}{#1}}
\newcommand{\red}[1]{\textcolor{ddlred}{#1}}
\newcommand{\green}[1]{\textcolor{ddlgreen}{#1}}

\newcommand{\RR}{\mathbb{R}}
\newcommand{\zero}{\bm{0}}
\newcommand{\bx}{\bm{x}}
\newcommand{\by}{\bm{y}}
\newcommand{\bz}{\bm{z}}
\newcommand{\bw}{\bm{w}}
\newcommand{\br}{\bm{r}}
\newcommand{\bq}{\bm{q}}
\newcommand{\bmu}{\bm{\mu}}
\newcommand{\blmbd}{\bm{\lambda}}
\newcommand{\calA}{\mathcal{A}}
\newcommand{\calL}{\mathcal{L}}
\newcommand{\calO}{\mathcal{O}}
\newcommand{\calS}{\mathcal{S}}
\newcommand{\prox}{\operatorname{prox}}
\newcommand{\dist}{\operatorname{dist}}
\newcommand{\st}{\textnormal{s.t.}}
\newcommand{\eps}{\varepsilon}

\setbeamertemplate{footline}{%
  \leavevmode%
  \hbox{%
    \begin{beamercolorbox}[wd=.23\paperwidth,ht=2.25ex,dp=1ex,center]{author in head/foot}%
      \usebeamerfont{author in head/foot}\insertshortauthor~(\insertshortinstitute)
    \end{beamercolorbox}%
    \begin{beamercolorbox}[wd=.54\paperwidth,ht=2.25ex,dp=1ex,center]{title in head/foot}%
      \usebeamerfont{title in head/foot}\insertshorttitle
    \end{beamercolorbox}%
    \begin{beamercolorbox}[wd=.23\paperwidth,ht=2.25ex,dp=1ex,right]{date in head/foot}%
      \usebeamerfont{date in head/foot}\insertshortdate{}\hspace*{1.4em}\insertframenumber{} / \inserttotalframenumber\hspace*{1ex}
    \end{beamercolorbox}%
  }%
  \vskip0pt%
}

\title[Moving-Center DDLALM]{Anytime Moving-Center DDLALM\\for Nonsmooth Composite Optimization}
\author[L. Li et al.]{Luxuan Li\inst{1}}
%  \and Kangkang Deng\inst{2} \and Chunfeng Cui\inst{1} \and Jiang Hu\inst{3}}
\institute[Beihang]{\inst{1} Beihang University}
%  \and \inst{2} National University of Defense Technology \and \inst{3} Tsinghua University}
\date[Research update]{Research update, August 2026}

% Insert a standalone divider slide before every section.
\AtBeginSection[]{
  \begin{frame}[plain,noframenumbering]
    \centering
    \vfill
    {\usebeamerfont{section title}\usebeamercolor[fg]{section title}%
      \Roman{section}.\enspace\insertsectionhead\par}
    \vfill
  \end{frame}
}

\begin{document}

\begin{frame}
  \titlepage
\end{frame}

\begin{frame}{Outline}
  \tableofcontents
\end{frame}

\section{Problem Setting}

\begin{frame}{Structured Nonsmooth Composite Optimization}
We consider
\begin{equation*}
  \min_{\bx\in\RR^n}\quad F(\bx):=f(\bx)+h\bigl(\calA(\bx)\bigr).
  \tag{P0}
\end{equation*}
\begin{itemize}
  \item $f:\RR^n\to\RR$ is continuously differentiable and may be nonconvex.
  \item $\calA:\RR^n\to\RR^m$ is continuously differentiable and may be nonlinear.
  \item $h:\RR^m\to\RR$ is convex, Lipschitz continuous, and possibly nonsmooth.
\end{itemize}
\end{frame}

\begin{frame}{Approximate KKT Point}
  \begin{itemize}
  \item Introduce $\by=\calA(\bx)$:
  \begin{equation}
    \label{Prob:split}
    \min_{\bx,\by}\quad f(\bx)+h(\by)
    \qquad \st\qquad \calA(\bx)-\by=\zero.   \tag{P}
  \end{equation}
  
  \item The split KKT system is
  \begin{equation*}
    \nabla f(\bx)+J_{\calA}(\bx)^\top\bq=\zero,
    \qquad \bq\in\partial h(\by),
    \qquad \calA(\bx)-\by=\zero.
  \end{equation*}

  \item \textbf{approximate $\eps$-KKT point.}
  Find $(\bx,\by,\bq)$ such that
  \begin{equation*}
    \|\nabla f(\bx)+J_{\calA}(\bx)^\top\bq\|\leq\eps,
    \quad \blue{\bq\in\partial h(\by)},
    \quad \|\calA(\bx)-\by\|\leq\eps.
  \end{equation*}
  We denote the maximum of these three residuals by
  $\mathcal R_{\rm split}(\bx,\by,\bq)$.

  \item The exact $\by$-proximal step makes the subdifferential condition \blue{exact}.
  \end{itemize}
\end{frame}

\begin{frame}{Main Result}
\small
The Candidate method uses the four-block order
\begin{equation*}
  \boxed{\bx\ \longrightarrow\ \blmbd\ \longrightarrow\ \bz\ \longrightarrow\ \bmu}
\end{equation*}
and the anytime schedule
\begin{equation*}
  r_k:=s_k^{-1}=\Theta(k^{-1/3}),
  \qquad
  \beta_k=\Theta(k^{-1/3}),
  \qquad
  \gamma_k=\Theta(k^{-1}).
\end{equation*}
Under the stated smoothness, lower-boundedness, and global Lipschitz
assumptions, one tail iterate satisfies
\begin{equation*}
  \mathcal R_{\rm split}
  (\bx^{\kappa+1},\by_\kappa^{\rm out},\bq_\kappa^{\rm out})
  =O(K^{-1/3}).
\end{equation*}
\begin{block}{Candidate guarantee}
Single loop, no algorithmic projection, no inner solve, fixed base penalty
$\rho$, two $\prox_{h/\rho}$ evaluations per iteration, and
$O(\eps^{-3})$ split-KKT complexity.
\end{block}
\end{frame}

\begin{frame}{Research Problems}
% 模仿

 Q1: Fixed penalty, if single-loop ?
  Q1: For single-loop first-order algorithms, if $\rho$ is fixed, to acclerate the rate?

  Q2: Can combinng the dual and smoothin improves the iterative complexity?
\end{frame}

\section{Reduced-Dual Reformulation}

\begin{frame}{Smoothed Augmented Lagrangian}

The smoothed augmented Lagrangian used in our method is
\[
{\mathcal P}(x,y,\lambda;z,\mu)
:=
\mathcal L_\rho(x,y,\lambda)
+\frac{r}{2}\|x-z\|^2
-\frac{1}{2s}\|\lambda-\mu\|^2.
\]
% We define the value function obtained by minimizing over the $\by$-block:

% \[
% F(\bx,\blmbd;\bz,\bmu)
% :=\min\limits_{\by}\mathcal P(\bx,\by,\blmbd;\bz,\bmu)
% \]
% \[
% \begin{aligned}
% F(\bx,\blmbd;\bz,\bmu)
% ={}&f(\bx)
% +\frac{r}{2}\|\bx-\bz\|^2
% -\frac{1}{2s}\|\blmbd-\bmu\|^2 \\
% &+\min\limits_{\by}
% \left\{
% h(\by)
% +\left\langle\blmbd,\calA(\bx)-\by\right\rangle
% +\frac{\rho}{2}\|\calA(\bx)-\by\|^2
% \right\}.
% \end{aligned}
% \]

\end{frame}


\begin{frame}{Value Function}
\small
1. From original problem to minmax problem;
2. Froem 
The value function
\begin{itemize}
  \item[\blue{$\triangleright$}]
  \textbf{Eliminate the $\by$-block through the $\rho$-prox:}
  \begin{equation*}
    d_s(\bx,\blmbd;\bmu)
    :=\min_{\by}\mathcal P_s(\bx,\by,\blmbd;\bmu),
    \qquad
    \red{\by(\blmbd;\bx)}
    =\operatorname{prox}_{h/\rho}
    \left(\calA(\bx)+\frac{\blmbd}{\rho}\right).
  \end{equation*}

  \item[\blue{$\triangleright$}]
  \textbf{Obtain an exact subdifferential certificate:}
  \begin{equation*}
    \blmbd+\rho\bigl[\calA(\bx)-\by(\blmbd;\bx)\bigr]
    \in\partial h\bigl(\by(\blmbd;\bx)\bigr).
  \end{equation*}

  \item[\blue{$\triangleright$}]
  \textbf{Apply Danskin's theorem:}
  \hfill
  \hyperlink{appendix:candidate-geometry}{%
    \beamergotobutton{Danskin derivation $\rightarrow$}%
  }
  \begin{equation*}
    \nabla_{\bx}d_s
    =\nabla f(\bx)+J_{\calA}(\bx)^\top
    \left\{\blmbd+\rho\bigl[\calA(\bx)-\by(\blmbd;\bx)\bigr]\right\},
  \end{equation*}
  \begin{equation*}
    \nabla_{\blmbd}d_s
    =\calA(\bx)-\by(\blmbd;\bx)-\frac{\blmbd-\bmu}{s},
    \qquad
    \nabla_{\bmu}d_s
    =\frac{\blmbd-\bmu}{s}.
  \end{equation*}
\end{itemize}


% \begin{block}{Moving-center identity}
% \begin{equation*}
%   \blue{\boxed{
%   \calA(\bx)-\by(\blmbd;\bx)
%   =\nabla_{\blmbd}d_s(\bx,\blmbd;\bmu)
%   +\nabla_{\bmu}d_s(\bx,\blmbd;\bmu)}}
% \end{equation*}
% The $\bmu$-gradient exactly complements the regularization term in the
% $\blmbd$-gradient.
% \end{block}

Under the strongly convex and strongly concave, it holds that:

\[
\boxed{
\max_\lambda d(\lambda,z,\mu)
=
p(z,\mu)
=
\min_x\max_\lambda\mathcal F(x,\lambda;z,\mu)
}
\]
\end{frame}

\section{Candidate Algorithm}

\begin{frame}{Candidate Algorithm: Four-Block Single Loop}
\footnotesize
Set $r_k:=1/s_k$ and use only the fixed-parameter mapping
\begin{equation*}
  \by(\blmbd;\bx)
  :=\prox_{h/\rho}\left(\calA(\bx)+\frac{\blmbd}{\rho}\right),
  \qquad
  \bq(\blmbd;\bx)
  :=\blmbd+\rho\bigl[\calA(\bx)-\by(\blmbd;\bx)\bigr].
\end{equation*}
\noindent\textit{Candidate notation: $\bq$ is the split-KKT multiplier and
$\bz$ is the primal center.}

Given $(\bx^k,\blmbd^k,\bz^k,\bmu^k)$, perform
\begin{equation*}
\begin{aligned}
  \by^k&=\by(\blmbd^k;\bx^k),
  \qquad \bq^k=\bq(\blmbd^k;\bx^k),\\
  \bx^{k+1}
  &=\bx^k-c\left[
  \nabla f(\bx^k)+J_{\calA}(\bx^k)^\top\bq^k
  +r_x(\bx^k-\bz^k)\right],\\
  \widehat\by^{k+1}
  &=\by(\blmbd^k;\bx^{k+1}),\\
  \blmbd^{k+1}
  &=\blmbd^k+\alpha\left[
  \calA(\bx^{k+1})-\widehat\by^{k+1}
  -r_k(\blmbd^k-\bmu^k)\right].
\end{aligned}
\end{equation*}
\begin{block}{Interpretation}
The first two blocks are one primal descent step and one explicit
reduced-dual gradient ascent step.
\end{block}
\end{frame}

\begin{frame}{Candidate Algorithm: Center Updates and Cost}
\footnotesize
Complete the iteration in the order
\begin{equation*}
  \boxed{\bx\ \longrightarrow\ \blmbd\ \longrightarrow\ \bz\ \longrightarrow\ \bmu}
\end{equation*}
with
\begin{equation*}
\begin{aligned}
  \bz^{k+1}
  &=\bz^k+\beta_k(\bx^{k+1}-\bz^k),\\
  \bmu^{k+1}
  &=\bmu^k+\gamma_k(\blmbd^{k+1}-\bmu^k).
\end{aligned}
\end{equation*}
Use the anytime schedules
\begin{equation*}
  r_k=c_r(k+k_0)^{-1/3},
  \qquad s_k=r_k^{-1},
\end{equation*}
\begin{equation*}
  \beta_k=\frac{\theta_z}{\ell_\rho}r_k,
  \qquad
  \gamma_k=\frac{\theta_\mu}{k+k_0},
  \qquad
  c,\alpha=\Theta(1/\ell_\rho).
\end{equation*}
Here $k_0\ge1$, $c_r>0$, $0<\theta_\mu\le1$, and $\theta_z>0$.
\begin{block}{Algorithmic properties}
No projection, no inner loop, and two evaluations of the same
$\prox_{h/\rho}$ mapping per iteration.
\end{block}
\end{frame}

\section{Complexity Result}

\begin{frame}{Candidate Assumptions and Step-Size Conditions}
\scriptsize
\begin{itemize}
  \setlength{\itemsep}{0.10cm}
  \item[(i)] $h$ is proper, closed, convex, finite-valued, and globally
  $L_h$-Lipschitz.
  \item[(ii)] $f$ and $\calA$ are continuously differentiable;
  $\nabla f$ and $J_{\calA}$ are globally Lipschitz, and
  $\|J_{\calA}(\bx)\|\le G_{\calA}$.
  \item[(iii)] $F_{\inf}:=\inf_{\bx}\{f(\bx)+h(\calA(\bx))\}>-\infty$.
\end{itemize}
Define
\begin{equation*}
  a_\rho:=L_f+L_{\calA}L_h+\rho G_{\calA}^2,
  \qquad
  \ell_\rho:=a_\rho+2G_{\calA}+\rho^{-1}.
\end{equation*}
Choose $r_x>3\ell_\rho$, $r_0\le\ell_\rho$, and fixed $c,\alpha>0$ such that
\begin{equation*}
  c<\frac{1}{r_x+\ell_\rho+r_0},
\end{equation*}
\begin{equation*}
  \alpha\le\min\left\{
  \frac{1}{11(\ell_\rho+r_0)},
  \frac{c^2m_0^2}{4(\ell_\rho+r_0)(1+cm_0)^2},
  \frac{1}{\rho^{-1}+r_0}
  \right\},
  \quad m_0:=r_x-\ell_\rho-r_0.
\end{equation*}
Choose $\theta_z>0$ sufficiently small so that every $\beta_k$ satisfies all
three caps stated in Lemma C4.2.  This choice is uniform in $k$, $K$, and
$\eps$; moreover $0<\gamma_k\le1$ follows from $k_0\ge1$ and
$0<\theta_\mu\le1$.
\end{frame}

\begin{frame}{Candidate Theorem: $O(\eps^{-3})$ Split-KKT Complexity}
\small
For each $k$, define the post-primal output
\begin{equation*}
  \by_k^{\rm out}:=\widehat\by^{k+1},
  \qquad
  \bq_k^{\rm out}:=\blmbd^k
  +\rho\bigl[\calA(\bx^{k+1})-\by_k^{\rm out}\bigr].
\end{equation*}
Under the preceding assumptions and parameter conditions, for every $K\ge4$
there exists one index
\begin{equation*}
  \kappa\in\{\lceil K/2\rceil,\ldots,K-1\}
\end{equation*}
such that
\begin{equation*}
  \bq_\kappa^{\rm out}\in\partial h(\by_\kappa^{\rm out}),
\end{equation*}
\begin{equation*}
  \left\|\nabla f(\bx^{\kappa+1})
  +J_{\calA}(\bx^{\kappa+1})^\top\bq_\kappa^{\rm out}\right\|
  \le C_{\rm stat}K^{-1/3},
\end{equation*}
\begin{equation*}
  \left\|\calA(\bx^{\kappa+1})-\by_\kappa^{\rm out}\right\|
  \le C_{\rm feas}K^{-1/3}.
\end{equation*}
\begin{block}{Complexity}
$K=O(\eps^{-3})$ iterations and $2O(\eps^{-3})$ fixed-parameter
$h$-prox evaluations; the result is tail-best-iterate.
\end{block}
\end{frame}

\section{Proof Sketch}

\begin{frame}{Proof Sketch I: Reduced Game and Frozen Descent}
\footnotesize
\begin{itemize}
  \setlength{\itemsep}{0.18cm}
  \item[\red{$\blacktriangleright$}]
  \textbf{Module 1: reduced geometry and bounded dual orbit.}
  Eliminating $\by$ gives a smooth nonconvex-concave payoff and an exact
  split-KKT multiplier:
  \begin{equation*}
    \nabla_{\blmbd}\mathcal D_\rho=\calA(\bx)-\by(\blmbd;\bx),
    \qquad
    \bq(\blmbd;\bx)\in\partial h(\by(\blmbd;\bx)).
  \end{equation*}
  The $\blmbd$- and $\bmu$-updates are convex combinations, hence the actual
  dual orbit is bounded without algorithmic projection.
  \hfill\hyperlink{appendix:candidate-geometry}{\beamergotobutton{Lemmas C1--C2}}

  \item[\red{$\blacktriangleright$}]
  \textbf{Module 2: frozen three-value descent.}
  Freeze $(r_k,\bmu^k)$ and identify the $(\bx,\blmbd,\bz)$ steps with
  Perturbed Smoothed GDA.  Its three-value Lyapunov satisfies
  \begin{equation*}
  \begin{aligned}
    \Delta\Psi_k\le{}&-c_x\|\Delta\bx^{k+1}\|^2
    -c_\lambda e_{\lambda,k}^2
    -c_z\beta_k\|\bz^k-\bx^{k+1}\|^2.
  \end{aligned}
  \end{equation*}
  \hfill\hyperlink{appendix:candidate-frozen}{\beamergotobutton{Lemmas C3--C4}}
\end{itemize}
\begin{block}{Core mechanism}
The three-value Lyapunov controls an ideal dual step directly and avoids the
$s_k^3\|\Delta\bx^{k+1}\|^2$ loss caused by exact-target tracking.
\end{block}
\end{frame}

\begin{frame}{Proof Sketch II: Moving Parameters and Rate Extraction}
\footnotesize
\begin{itemize}
  \setlength{\itemsep}{0.16cm}
  \item[\red{$\blacktriangleright$}]
  \textbf{Module 3: summable transport and master budget.}
  Moving $\bmu^k$ and varying $r_k$ contribute only
  $O(r_k\gamma_k)$ and $O(|r_{k+1}-r_k|)$.  Both sequences are summable, so
  frozen descent telescopes to
  \begin{equation*}
    \sum_{k=0}^{\infty}\left[
    \|\Delta\bx^{k+1}\|^2+e_{\lambda,k}^2
    +\beta_k\|\bz^k-\bx^{k+1}\|^2\right]<\infty.
  \end{equation*}
  \hfill\hyperlink{appendix:candidate-transport}{\beamergotobutton{Lemmas C5--C6}}

  \item[\red{$\blacktriangleright$}]
  \textbf{Module 4: split-KKT reconstruction and tail balance.}
  At one common $\kappa\in\{\lceil K/2\rceil,\ldots,K-1\}$,
  \begin{equation*}
    \|\Delta\bx^{\kappa+1}\|+e_{\lambda,\kappa}=O(K^{-1/2}),
    \qquad
    \|\bz^\kappa-\bx^{\kappa+1}\|
    +r_\kappa\|\blmbd^\kappa-\bmu^\kappa\|=O(K^{-1/3}).
  \end{equation*}
  The update identities convert these quantities into stationarity and
  feasibility residuals of order $O(K^{-1/3})$.
  \hfill\hyperlink{appendix:candidate-residual}{\beamergotobutton{Lemmas C7--C8}}
\end{itemize}
\begin{block}{Final rate}
The exact subgradient residual is zero; the other two split-KKT residuals are
$O(K^{-1/3})$.  Therefore $K=O(\eps^{-3})$.
\end{block}
\end{frame}

\begin{frame}{Candidate Proof Status and Scope}
\small
\begin{itemize}
  \setlength{\itemsep}{0.18cm}
  \item The four-block algorithm on the slides is exactly the algorithm used
  in the varying-$s_k$ proof.
  \item The frozen-center step is the translated fixed-perturbation instance
  of Li--Nagarajan--Pan--Zhang's Perturbed Smoothed GDA
  three-value analysis~{\normalfont[4]}.
  \item The moving-$\bmu$ transport, varying-$r_k$ drift, no-projection
  boundedness, and split-KKT reconstruction are proved separately in
  the complete-proof appendix.
  \item The theorem requires globally Lipschitz $h$ and gives a
  deterministic tail-best-iterate split-KKT result.
  \item The result does not establish the stronger constant-$s$ route and
  does not claim last-iterate complexity.
\end{itemize}
\end{frame}

\section{References}

\begin{frame}{References}
\footnotesize
\begin{enumerate}
  \item[{[1]}] H. Zhang, N. Marumo, T. K. Pong, and A. Takeda,
  ``Complexity and Convergence Analysis of a Single-Loop SDCAM for
  Lipschitz Composite Optimization and Beyond,'' arXiv:2512.24059, 2025.

  \item[{[2]}] K. Sun and X. A. Sun,
  ``Dual Descent Augmented Lagrangian Method and Alternating Direction
  Method of Multipliers,'' \emph{SIAM Journal on Optimization},
  34(2):1679--1707, 2024.

  \item[{[3]}] J. M. Danskin,
  ``The Theory of Max--Min, with Applications,''
  \emph{SIAM Journal on Applied Mathematics},
  14(4):641--664, 1966.

  \item[{[4]}] J. Li, M. Nagarajan, S. Pan, and N. Zhang,
  ``Smoothing Meets Perturbation: Unified and Tight Analysis for
  Nonconvex--Concave Minimax Optimization,'' arXiv:2602.14185v2, 2026.
\end{enumerate}
\end{frame}

\appendix

\begin{frame}{Appendix: Complete Proofs}
\hypertarget{appendix:contents}{}
\small
\begin{enumerate}
  \item \hyperlink{appendix:candidate-geometry}{Lemmas C1--C2: reduced geometry and bounded dual orbit}
  \item \hyperlink{appendix:candidate-frozen}{Lemmas C3--C4: three-value Lyapunov and frozen descent}
  \item \hyperlink{appendix:candidate-transport}{Lemmas C5--C6: moving-parameter transport and master budget}
  \item \hyperlink{appendix:candidate-residual}{Lemmas C7--C8: residual reconstruction and tail extraction}
  \item \hyperlink{appendix:candidate-main-theorem}{Main Theorem: complete complexity proof}
\end{enumerate}
\begin{block}{Detail level}
Every appendix result includes its assumptions, full derivation, and the exact point where it enters the main theorem.
\end{block}
\end{frame}

\begin{frame}{Appendix: Lemma C1.1---Reduced-Payoff Calculus
\hfill\hyperlink{appendix:contents}{\beamerreturnbutton{Contents}}}
\hypertarget{appendix:candidate-geometry}{}
\scriptsize
Define
\begin{equation*}
  \by(\blmbd;\bx):=\prox_{h/\rho}
  \left(\calA(\bx)+\frac{\blmbd}{\rho}\right),
  \qquad
  \bq(\blmbd;\bx):=\blmbd+\rho[\calA(\bx)-\by(\blmbd;\bx)],
\end{equation*}
and
\begin{equation*}
  \mathcal D_\rho(\bx,\blmbd)
  :=\min_{\by}\left\{f(\bx)+h(\by)
  +\langle\blmbd,\calA(\bx)-\by\rangle
  +\frac{\rho}{2}\|\calA(\bx)-\by\|^2\right\}.
\end{equation*}
Then
\begin{equation*}
  \nabla_{\bx}\mathcal D_\rho
  =\nabla f(\bx)+J_{\calA}(\bx)^\top\bq(\blmbd;\bx),
  \qquad
  \nabla_{\blmbd}\mathcal D_\rho
  =\calA(\bx)-\by(\blmbd;\bx),
\end{equation*}
and
\begin{equation*}
  \bq(\blmbd;\bx)\in\partial h(\by(\blmbd;\bx)),
  \qquad \|\bq(\blmbd;\bx)\|\le L_h.
\end{equation*}
\noindent\textbf{Proof.}
Completing the square gives
\begin{equation*}
  h(\by)+\frac{\rho}{2}
  \left\|\by-\calA(\bx)-\frac{\blmbd}{\rho}\right\|^2
  -\frac{1}{2\rho}\|\blmbd\|^2.
\end{equation*}
Thus the minimizer is unique.  Its optimality condition is
\begin{equation*}
  0\in\partial h(\by)+\rho[\by-\calA(\bx)]-\blmbd,
\end{equation*}
which gives the exact certificate.  Global Lipschitz continuity of $h$ gives
$\|\bq\|\le L_h$.  Danskin's theorem applied to the unique minimizer gives
the two displayed gradients.
\end{frame}

\begin{frame}{Appendix: Lemma C1.2---Joint Smoothness
\hfill\hyperlink{appendix:contents}{\beamerreturnbutton{Contents}}}
\scriptsize
Moreau decomposition yields
\begin{equation*}
  \bq(\blmbd;\bx)=\prox_{\rho h^*}
  (\blmbd+\rho\calA(\bx)).
\end{equation*}
Hence nonexpansiveness gives
\begin{equation*}
  \|\bq(\blmbd;\bx)-\bq(\blmbd';\bx')\|
  \le \|\blmbd-\blmbd'\|+\rho G_{\calA}\|\bx-\bx'\|.
\end{equation*}
Using $\|\bq\|\le L_h$, $\|J_{\calA}\|\le G_{\calA}$, and the Lipschitz
constants of $\nabla f$ and $J_{\calA}$,
\begin{equation*}
\begin{aligned}
  &\|\nabla_{\bx}\mathcal D_\rho(\bx,\blmbd)
  -\nabla_{\bx}\mathcal D_\rho(\bx',\blmbd')\|\\
  &\quad\le
  (L_f+L_{\calA}L_h+\rho G_{\calA}^2)\|\bx-\bx'\|
  +G_{\calA}\|\blmbd-\blmbd'\|.
\end{aligned}
\end{equation*}
Moreover,
\begin{equation*}
  \calA(\bx)-\by(\blmbd;\bx)
  =\rho^{-1}[\bq(\blmbd;\bx)-\blmbd].
\end{equation*}
Firm nonexpansiveness makes $I-\prox_{\rho h^*}$ nonexpansive.  Therefore
\begin{equation*}
  \|\nabla_{\blmbd}\mathcal D_\rho(\bx,\blmbd)
  -\nabla_{\blmbd}\mathcal D_\rho(\bx',\blmbd')\|
  \le G_{\calA}\|\bx-\bx'\|+\rho^{-1}\|\blmbd-\blmbd'\|.
\end{equation*}
Thus $\mathcal D_\rho$ is jointly $\ell_\rho$-smooth for
$\ell_\rho:=L_f+L_{\calA}L_h+\rho G_{\calA}^2+2G_{\calA}+\rho^{-1}$.
\end{frame}

\begin{frame}{Appendix: Lemma C1.3---Concavity and Exact Maximum
\hfill\hyperlink{appendix:contents}{\beamerreturnbutton{Contents}}}
\scriptsize
For fixed $\bx$, $\mathcal D_\rho(\bx,\cdot)$ is the pointwise infimum of
affine functions of $\blmbd$, hence it is concave.  It also satisfies
\begin{equation*}
  \max_{\blmbd\in\RR^m}\mathcal D_\rho(\bx,\blmbd)
  =f(\bx)+h(\calA(\bx)).
\end{equation*}
\noindent\textbf{Proof.}
Taking $\by=\calA(\bx)$ in the inner minimization gives, for every $\blmbd$,
\begin{equation*}
  \mathcal D_\rho(\bx,\blmbd)
  \le f(\bx)+h(\calA(\bx)).
\end{equation*}
Because $h$ is finite-valued and convex, choose
$\bar\blmbd\in\partial h(\calA(\bx))$.  Then $\by=\calA(\bx)$ satisfies
\begin{equation*}
  0\in\partial h(\by)-\bar\blmbd
  +\rho[\by-\calA(\bx)],
\end{equation*}
so it is the inner minimizer and equality holds at $\bar\blmbd$.  Global
Lipschitz continuity gives $\|\bar\blmbd\|\le L_h$.  Consequently, every
maximization needed later may select a maximizer in the $L_h$-ball.
\begin{block}{Use in the proof}
This exact maximum identity supplies the uniform lower bound for the
three-value Lyapunov function in Lemma C6.
\end{block}
\end{frame}

\begin{frame}{Appendix: Lemma C2---No-Projection Boundedness
\hfill\hyperlink{appendix:contents}{\beamerreturnbutton{Contents}}}
\hypertarget{appendix:candidate-bounded}{}
\scriptsize
Let $M:=\max\{\|\blmbd^0\|,\|\bmu^0\|,L_h\}$.  If
$\alpha\le(\rho^{-1}+r_0)^{-1}$ and $0\le\gamma_k\le1$, then
\begin{equation*}
  \|\blmbd^k\|\le M,
  \qquad \|\bmu^k\|\le M,
  \qquad k\ge0.
\end{equation*}
\noindent\textbf{Proof.}
Set
\begin{equation*}
  \widehat\bq^{k+1}:=\blmbd^k
  +\rho[\calA(\bx^{k+1})-\widehat\by^{k+1}].
\end{equation*}
Lemma C1 gives $\|\widehat\bq^{k+1}\|\le L_h$.  Substitution into the
$\blmbd$-update yields
\begin{equation*}
  \blmbd^{k+1}
  =[1-\alpha(\rho^{-1}+r_k)]\blmbd^k
  +\frac{\alpha}{\rho}\widehat\bq^{k+1}
  +\alpha r_k\bmu^k.
\end{equation*}
Since $0<r_k\le r_0$, the three coefficients are nonnegative and sum to one.
Thus, whenever $\|\blmbd^k\|,\|\bmu^k\|\le M$,
\begin{equation*}
  \|\blmbd^{k+1}\|\le M.
\end{equation*}
The update
$\bmu^{k+1}=(1-\gamma_k)\bmu^k+\gamma_k\blmbd^{k+1}$ then gives
$\|\bmu^{k+1}\|\le M$.  The base case follows from the definition of $M$,
so induction proves the claim.  Define $B_M:=\{v:\|v\|\le M\}$.  Choose $R>M$ and
$Y:=\{\blmbd:\|\blmbd\|\le R\}$.  The actual iterates lie strictly inside
$Y$; projection onto $Y$ is used only in the proof.
\end{frame}

\begin{frame}{Appendix: Lemma C3.1---Three Value Functions
\hfill\hyperlink{appendix:contents}{\beamerreturnbutton{Contents}}}
\hypertarget{appendix:candidate-frozen}{}
\scriptsize
For fixed $r\in(0,r_0]$, $\bmu\in B_M$, and $\bz\in\RR^n$, define
\begin{equation*}
  \mathcal F_{r,\bmu}(\bx,\blmbd,\bz)
  :=\mathcal D_\rho(\bx,\blmbd)
  +\frac{r_x}{2}\|\bx-\bz\|^2
  -\frac r2\|\blmbd-\bmu\|^2,
\end{equation*}
\begin{equation*}
  d_{r,\bmu}(\blmbd,\bz)
  :=\min_{\bx}\mathcal F_{r,\bmu}(\bx,\blmbd,\bz),
  \qquad
  p_{r,\bmu}(\bz)
  :=\min_{\bx}\max_{\blmbd\in Y}\mathcal F_{r,\bmu}(\bx,\blmbd,\bz).
\end{equation*}
The minima exist, and the minimizer defining $d_{r,\bmu}$ is unique.
\noindent\textbf{Proof.}
Joint $\ell_\rho$-smoothness implies that
$\mathcal D_\rho(\cdot,\blmbd)$ is $\ell_\rho$-weakly convex.  Therefore
\begin{equation*}
  \mathcal F_{r,\bmu}(\cdot,\blmbd,\bz)
  \text{ is }(r_x-\ell_\rho)\text{-strongly convex}.
\end{equation*}
It is finite and coercive, so its unique minimizer exists.  For each
$\blmbd\in Y$, the same strong-convexity modulus holds.  The pointwise maximum
over compact $Y$ preserves strong convexity, finiteness, and coercivity.
Hence the minimum defining $p_{r,\bmu}$ also exists.
\end{frame}

\begin{frame}{Appendix: Lemma C3.2---Three-Value Lyapunov
\hfill\hyperlink{appendix:contents}{\beamerreturnbutton{Contents}}}
\scriptsize
Define
\begin{equation*}
  \Psi_{r,\bmu}(\bx,\blmbd,\bz)
  :=\mathcal F_{r,\bmu}(\bx,\blmbd,\bz)
  -2d_{r,\bmu}(\blmbd,\bz)+2p_{r,\bmu}(\bz).
\end{equation*}
Its exact decomposition is
\begin{equation*}
  \Psi_{r,\bmu}
  =(\mathcal F_{r,\bmu}-d_{r,\bmu})
  +(p_{r,\bmu}-d_{r,\bmu})+p_{r,\bmu}.
\end{equation*}
The first two bracketed terms are nonnegative because $d$ is a pointwise
minimum and
\begin{equation*}
  p_{r,\bmu}(\bz)
  =\min_{\bx}\max_{\blmbd'\in Y}\mathcal F_{r,\bmu}(\bx,\blmbd',\bz)
  \ge d_{r,\bmu}(\blmbd,\bz).
\end{equation*}
Let
\begin{equation*}
  \bx_{r,\bmu}(\blmbd,\bz)
  :=\operatorname*{arg\,min}_{\bx}\mathcal F_{r,\bmu}(\bx,\blmbd,\bz)
\end{equation*}
and define the ideal dual step
\begin{equation*}
  \blmbd_+(\bz):=\operatorname{proj}_Y\left[
  \blmbd+\alpha\nabla_{\blmbd}\mathcal F_{r,\bmu}
  (\bx_{r,\bmu}(\blmbd,\bz),\blmbd,\bz)\right].
\end{equation*}
The three terms of $\Psi$ respectively measure primal suboptimality, dual
ascent error, and the proximal-center value.
\end{frame}

\begin{frame}{Appendix: Lemma C4.1---Frozen-Center Transfer
\hfill\hyperlink{appendix:contents}{\beamerreturnbutton{Contents}}}
\scriptsize
Fix $(r,\bmu)$ and set
\begin{equation*}
  \boldsymbol\nu:=\blmbd-\bmu,
  \qquad
  \widetilde D_{\bmu}(\bx,\boldsymbol\nu)
  :=\mathcal D_\rho(\bx,\boldsymbol\nu+\bmu).
\end{equation*}
Then $\widetilde D_{\bmu}$ is $\ell_\rho$-smooth and concave in
$\boldsymbol\nu$.  The translated dual set $Y-\bmu$ is compact and convex,
contains the origin because $\bmu\in B_M\subset Y$, and has the same diameter
as $Y$.  Moreover,
\begin{equation*}
  \mathcal F_{r,\bmu}
  =\widetilde D_{\bmu}(\bx,\boldsymbol\nu)
  +\frac{r_x}{2}\|\bx-\bz\|^2-\frac r2\|\boldsymbol\nu\|^2.
\end{equation*}
Thus the frozen $(\bx,\blmbd,\bz)$ update is exactly Perturbed Smoothed GDA
on the translated problem.  Li--Nagarajan--Pan--Zhang~{\normalfont[4]},
Lemmas 4.4--4.5, give
\begin{equation*}
\begin{aligned}
  \Psi^+-\Psi\le{}&-\frac1{8c}\|\bx^+-\bx\|^2
  -\frac1{8\alpha}\|\blmbd-\blmbd_+(\bz)\|^2
  -\frac{r_x\beta}{8}\|\bz-\bx^+\|^2\\
  &+24r_x\beta
  \|\bx_r^*(\bz)-\bx_r(\blmbd_+(\bz),\bz)\|^2,
\end{aligned}
\end{equation*}
and
\begin{equation*}
  \|\bx_r^*(\bz)-\bx_r(\blmbd_+(\bz),\bz)\|
  \le\Omega(r)\|\blmbd-\blmbd_+(\bz)\|.
\end{equation*}
The perturbation-change term in their Lemma 4.4 is zero because $r$ is frozen.
\end{frame}

\begin{frame}{Appendix: Lemma C4.2---Parameter Map and Absorption
\hfill\hyperlink{appendix:contents}{\beamerreturnbutton{Contents}}}
\scriptsize
The parameter map to Condition 3.2 of~{\normalfont[4]} is
\begin{equation*}
  (\ell,r_x,c,\alpha,\beta,r_y)
  \longmapsto(\ell_\rho,r_x,c,\alpha,\beta,r).
\end{equation*}
Its dual error coefficient becomes
\begin{equation*}
  \Omega(r)
  :=\frac1{\sqrt r}
  \frac{(r_x-\ell_\rho-r)
  +\alpha(\ell_\rho+r)(3r_x-2\ell_\rho-2r)}
  {\alpha(r_x-\ell_\rho-r)^{3/2}}.
\end{equation*}
Assume $r_x>3\ell_\rho$, $0<r\le r_0\le\ell_\rho$, the stated bounds on
$c$ and $\alpha$, and
\begin{equation*}
  0<\beta\le\min\left\{
  \frac1{36},
  \frac{(r_x-\ell_\rho-r)^2}
  {384r_x(r_x+\ell_\rho+r)^2},
  \frac1{384r_x\alpha\Omega(r)^2}
  \right\}.
\end{equation*}
Then
\begin{equation*}
  24r_x\beta\Omega(r)^2\le\frac1{16\alpha}.
\end{equation*}
Substitution in Lemma C4.1 gives the complete frozen descent inequality
\begin{equation*}
\begin{aligned}
  \Psi^+-\Psi\le{}&-\frac1{8c}\|\bx^+-\bx\|^2
  -\frac1{16\alpha}\|\blmbd-\blmbd_+(\bz)\|^2
  -\frac{r_x\beta}{8}\|\bz-\bx^+\|^2.
\end{aligned}
\end{equation*}
Since $\Omega(r)\le C_\Omega/\sqrt r$, choosing
$\beta_k=(\theta_z/\ell_\rho)r_k$ with sufficiently small fixed $\theta_z$
makes these bounds uniform in $k$.
\end{frame}

\begin{frame}{Appendix: Lemma C5.1---Moving-Center Transport
\hfill\hyperlink{appendix:contents}{\beamerreturnbutton{Contents}}}
\hypertarget{appendix:candidate-transport}{}
\scriptsize
For $\|\bmu\|,\|\bmu'\|\le M$, $\blmbd\in Y$, and
$u=(\bx,\blmbd,\bz)$,
\begin{equation*}
\begin{aligned}
  |\Psi_{r,\bmu'}(u)-\Psi_{r,\bmu}(u)|
  \le{}&5r(R+M)\|\bmu'-\bmu\|\\
  &+\frac{5r}{2}\|\bmu'-\bmu\|^2.
\end{aligned}
\end{equation*}
\noindent\textbf{Proof.}
Let $\delta\bmu:=\bmu'-\bmu$.  The only center-dependent term satisfies
\begin{equation*}
  -\frac r2\|\blmbd-\bmu'\|^2
  +\frac r2\|\blmbd-\bmu\|^2
  =r\langle\blmbd-\bmu,\delta\bmu\rangle
  -\frac r2\|\delta\bmu\|^2.
\end{equation*}
Because $\|\blmbd\|\le R$ and $\|\bmu\|\le M$,
\begin{equation*}
  \sup_{\bx,\blmbd,\bz}
  |\mathcal F_{r,\bmu'}-\mathcal F_{r,\bmu}|
  \le r(R+M)\|\delta\bmu\|+\frac r2\|\delta\bmu\|^2=:\delta.
\end{equation*}
Uniform perturbation by $\delta$ changes each minimum, maximum, and min--max
value by at most $\delta$.  Since
$\Psi=\mathcal F-2d+2p$, the total change is at most $5\delta$.
For the algorithm,
\begin{equation*}
  \|\bmu^{k+1}-\bmu^k\|
  =\gamma_k\|\blmbd^{k+1}-\bmu^k\|
  \le2M\gamma_k.
\end{equation*}
Using $\gamma_k\le1$, the transport is bounded by $C_\mu r_k\gamma_k$.
\end{frame}

\begin{frame}{Appendix: Lemma C5.2---Perturbation-Scale Drift
\hfill\hyperlink{appendix:contents}{\beamerreturnbutton{Contents}}}
\scriptsize
For fixed $\bmu$ and any $r,r'>0$,
\begin{equation*}
  |\Psi_{r',\bmu}(u)-\Psi_{r,\bmu}(u)|
  \le C_r|r'-r|,
  \qquad C_r:=\frac52(R+M)^2.
\end{equation*}
\noindent\textbf{Proof.}
The surrogate difference is exactly
\begin{equation*}
  \mathcal F_{r',\bmu}(u)-\mathcal F_{r,\bmu}(u)
  =-\frac{r'-r}{2}\|\blmbd-\bmu\|^2.
\end{equation*}
Therefore its uniform magnitude is at most
$\frac12(R+M)^2|r'-r|$.  Applying the same value-perturbation argument as in
Lemma C5.1 and the coefficient sum $1+2+2=5$ proves the claim.

For
\begin{equation*}
  r_k=c_r(k+k_0)^{-1/3},
  \qquad
  \gamma_k=\frac{\theta_\mu}{k+k_0},
\end{equation*}
$r_k$ is decreasing to zero, hence
\begin{equation*}
  \sum_{k=0}^{\infty}|r_{k+1}-r_k|=r_0.
\end{equation*}
Also
\begin{equation*}
  r_k\gamma_k=c_r\theta_\mu(k+k_0)^{-4/3},
  \qquad
  \sum_{k=0}^{\infty}r_k\gamma_k<\infty.
\end{equation*}
Thus both moving-parameter errors are summable with constants independent of
$K$ and $\eps$.
\end{frame}

\begin{frame}{Appendix: Lemma C6.1---Master One-Step Inequality
\hfill\hyperlink{appendix:contents}{\beamerreturnbutton{Contents}}}
\hypertarget{appendix:candidate-budget}{}
\scriptsize
Let
\begin{equation*}
  u^k:=(\bx^k,\blmbd^k,\bz^k),
  \quad V_k:=\Psi_{r_k,\bmu^k}(u^k),
  \quad e_{\lambda,k}:=\|\blmbd^k-\blmbd_{+,k}(\bz^k)\|.
\end{equation*}
Insert two intermediate values:
\begin{equation*}
\begin{aligned}
  V_{k+1}-V_k
  ={}&[\Psi_{r_k,\bmu^k}(u^{k+1})-\Psi_{r_k,\bmu^k}(u^k)]\\
  &+[\Psi_{r_k,\bmu^{k+1}}(u^{k+1})
  -\Psi_{r_k,\bmu^k}(u^{k+1})]\\
  &+[\Psi_{r_{k+1},\bmu^{k+1}}(u^{k+1})
  -\Psi_{r_k,\bmu^{k+1}}(u^{k+1})].
\end{aligned}
\end{equation*}
Lemma C4 controls the first bracket.  Lemma C5.1 controls the second, and
Lemma C5.2 controls the third.  Hence
\begin{equation*}
\begin{aligned}
  V_{k+1}-V_k\le{}&-\frac1{8c}\|\Delta\bx^{k+1}\|^2
  -\frac1{16\alpha}e_{\lambda,k}^2
  -\frac{r_x\beta_k}{8}\|\bz^k-\bx^{k+1}\|^2\\
  &+C_r|r_{k+1}-r_k|+C_\mu r_k\gamma_k.
\end{aligned}
\end{equation*}
Every positive error on the right is summable by Lemma C5.2.
\end{frame}

\begin{frame}{Appendix: Lemma C6.2---Uniform Lower Bound
\hfill\hyperlink{appendix:contents}{\beamerreturnbutton{Contents}}}
\scriptsize
For every $k$,
\begin{equation*}
  V_k\ge\underline V
  :=F_{\inf}-\frac{r_0}{2}(L_h+M)^2.
\end{equation*}
\noindent\textbf{Proof.}
Lemma C3 gives
\begin{equation*}
  \Psi_{r,\bmu}
  =(\mathcal F-d)+(p-d)+p\ge p.
\end{equation*}
For each $\bx$, Lemma C1.3 provides
$\bar\blmbd(\bx)\in\partial h(\calA(\bx))\subseteq Y$ such that
\begin{equation*}
  \mathcal D_\rho(\bx,\bar\blmbd(\bx))
  =f(\bx)+h(\calA(\bx)),
  \qquad \|\bar\blmbd(\bx)\|\le L_h.
\end{equation*}
Therefore
\begin{equation*}
\begin{aligned}
  \max_{\blmbd\in Y}\mathcal F_{r,\bmu}(\bx,\blmbd,\bz)
  &\ge f(\bx)+h(\calA(\bx))
  +\frac{r_x}{2}\|\bx-\bz\|^2
  -\frac r2\|\bar\blmbd(\bx)-\bmu\|^2\\
  &\ge F_{\inf}-\frac{r_0}{2}(L_h+M)^2.
\end{aligned}
\end{equation*}
Taking the minimum over $\bx$ proves the same lower bound for $p$, and hence
for $V_k$.
\end{frame}

\begin{frame}{Appendix: Lemma C6.3---Finite Master Budget
\hfill\hyperlink{appendix:contents}{\beamerreturnbutton{Contents}}}
\scriptsize
Sum Lemma C6.1 from $k=0$ to $N-1$.  After telescoping and applying
Lemma C6.2,
\begin{equation*}
\begin{aligned}
  &\sum_{k=0}^{N-1}\left[
  \frac1{8c}\|\Delta\bx^{k+1}\|^2
  +\frac1{16\alpha}e_{\lambda,k}^2
  +\frac{r_x\beta_k}{8}\|\bz^k-\bx^{k+1}\|^2\right]\\
  &\quad\le V_0-\underline V
  +C_r\sum_{k=0}^{\infty}|r_{k+1}-r_k|
  +C_\mu\sum_{k=0}^{\infty}r_k\gamma_k.
\end{aligned}
\end{equation*}
The right-hand side is a finite constant independent of $N$, $K$, and
$\eps$.  Letting $N\to\infty$ and absorbing the fixed positive coefficients
gives a constant $B<\infty$ such that
\begin{equation*}
  \sum_{k=0}^{\infty}\left[
  \|\Delta\bx^{k+1}\|^2+e_{\lambda,k}^2
  +\beta_k\|\bz^k-\bx^{k+1}\|^2\right]\le B.
\end{equation*}
\begin{block}{Use in the proof}
This horizon-independent budget is the sole averaging input in the tail-index
argument.  No last-iterate claim is used.
\end{block}
\end{frame}

\begin{frame}{Appendix: Lemma C7.1---Primal Target Error
\hfill\hyperlink{appendix:contents}{\beamerreturnbutton{Contents}}}
\hypertarget{appendix:candidate-residual}{}
\scriptsize
Let
\begin{equation*}
  \bx_k^*:=\operatorname*{arg\,min}_{\bx}
  \mathcal F_{r_k,\bmu^k}(\bx,\blmbd^k,\bz^k),
  \qquad m_x:=r_x-\ell_\rho>0.
\end{equation*}
Then
\begin{equation*}
  \|\bx^{k+1}-\bx_k^*\|
  \le\left(1+\frac1{cm_x}\right)\|\Delta\bx^{k+1}\|.
\end{equation*}
\noindent\textbf{Proof.}
The $\bx$-subproblem is $m_x$-strongly convex and
$\nabla_{\bx}\mathcal F_{r_k,\bmu^k}(\bx_k^*,\blmbd^k,\bz^k)=0$.
Strong monotonicity gives
\begin{equation*}
  m_x\|\bx^k-\bx_k^*\|
  \le\|\nabla_{\bx}\mathcal F_{r_k,\bmu^k}
  (\bx^k,\blmbd^k,\bz^k)\|.
\end{equation*}
The algorithmic $\bx$-step yields
\begin{equation*}
  \nabla_{\bx}\mathcal F_{r_k,\bmu^k}
  (\bx^k,\blmbd^k,\bz^k)
  =-c^{-1}\Delta\bx^{k+1}.
\end{equation*}
Thus
$\|\bx^k-\bx_k^*\|\le(cm_x)^{-1}\|\Delta\bx^{k+1}\|$.
The triangle inequality with
$\bx^{k+1}-\bx_k^*=(\bx^{k+1}-\bx^k)+(\bx^k-\bx_k^*)$
proves the claim.
\end{frame}

\begin{frame}{Appendix: Lemma C7.2---Actual Dual Increment
\hfill\hyperlink{appendix:contents}{\beamerreturnbutton{Contents}}}
\scriptsize
Because Lemma C2 keeps the raw dual update in $B_M\subset Y$,
\begin{equation*}
  \blmbd^{k+1}=\operatorname{proj}_Y\left[
  \blmbd^k+\alpha\nabla_{\blmbd}\mathcal F_{r_k,\bmu^k}
  (\bx^{k+1},\blmbd^k,\bz^k)\right].
\end{equation*}
The ideal step is the same expression with $\bx^{k+1}$ replaced by
$\bx_k^*$.  Projection nonexpansiveness and the cross-Lipschitz bound in
Lemma C1 give
\begin{equation*}
\begin{aligned}
  \|\blmbd^{k+1}-\blmbd_{+,k}(\bz^k)\|
  &\le\alpha G_{\calA}\|\bx^{k+1}-\bx_k^*\|\\
  &\le\alpha G_{\calA}
  \left(1+\frac1{cm_x}\right)\|\Delta\bx^{k+1}\|.
\end{aligned}
\end{equation*}
Adding and subtracting $\blmbd_{+,k}(\bz^k)$ yields
\begin{equation*}
  \|\blmbd^{k+1}-\blmbd^k\|
  \le e_{\lambda,k}+C_{\lambda x}\|\Delta\bx^{k+1}\|,
\end{equation*}
where
\begin{equation*}
  C_{\lambda x}:=\alpha G_{\calA}
  \left(1+\frac1{cm_x}\right).
\end{equation*}
This proves that the ideal dual error controlled by the Lyapunov budget also
controls the actual multiplier increment.
\end{frame}

\begin{frame}{Appendix: Lemma C8.1---Exact Certificate and Stationarity
\hfill\hyperlink{appendix:contents}{\beamerreturnbutton{Contents}}}
\scriptsize
Define
\begin{equation*}
  \by_k^{\rm out}:=\widehat\by^{k+1},
  \qquad
  \bq_k^{\rm out}:=\blmbd^k
  +\rho[\calA(\bx^{k+1})-\by_k^{\rm out}].
\end{equation*}
Lemma C1 gives the exact relation
\begin{equation*}
  \bq_k^{\rm out}\in\partial h(\by_k^{\rm out}).
\end{equation*}
The $\bx$-step is equivalent to
\begin{equation*}
  \nabla_{\bx}\mathcal D_\rho(\bx^k,\blmbd^k)
  +r_x(\bx^k-\bz^k)=-c^{-1}\Delta\bx^{k+1}.
\end{equation*}
Since
\begin{equation*}
  \nabla_{\bx}\mathcal D_\rho(\bx^{k+1},\blmbd^k)
  =\nabla f(\bx^{k+1})
  +J_{\calA}(\bx^{k+1})^\top\bq_k^{\rm out},
\end{equation*}
adding and subtracting the gradient at $\bx^k$, applying the fixed-$\blmbd$
$a_\rho$-smoothness bound, and using
$\bx^k-\bz^k=(\bx^{k+1}-\bz^k)-\Delta\bx^{k+1}$ give
\begin{equation*}
\begin{aligned}
  &\|\nabla f(\bx^{k+1})
  +J_{\calA}(\bx^{k+1})^\top\bq_k^{\rm out}\|\\
  &\quad\le(c^{-1}+r_x+a_\rho)\|\Delta\bx^{k+1}\|
  +r_x\|\bx^{k+1}-\bz^k\|.
\end{aligned}
\end{equation*}
\end{frame}

\begin{frame}{Appendix: Lemma C8.2---Feasibility Reconstruction
\hfill\hyperlink{appendix:contents}{\beamerreturnbutton{Contents}}}
\scriptsize
The dual update gives the exact identity
\begin{equation*}
  \frac{\blmbd^{k+1}-\blmbd^k}{\alpha}
  =\calA(\bx^{k+1})-\by_k^{\rm out}
  -r_k(\blmbd^k-\bmu^k).
\end{equation*}
Rearranging and using Lemma C2,
\begin{equation*}
\begin{aligned}
  \|\calA(\bx^{k+1})-\by_k^{\rm out}\|
  &\le\alpha^{-1}\|\blmbd^{k+1}-\blmbd^k\|
  +r_k\|\blmbd^k-\bmu^k\|\\
  &\le\alpha^{-1}\|\blmbd^{k+1}-\blmbd^k\|+2Mr_k.
\end{aligned}
\end{equation*}
Combining this inequality with Lemma C7.2 gives
\begin{equation*}
\begin{aligned}
  \|\calA(\bx^{k+1})-\by_k^{\rm out}\|
  \le{}&\alpha^{-1}e_{\lambda,k}
  +\alpha^{-1}C_{\lambda x}\|\Delta\bx^{k+1}\|+2Mr_k.
\end{aligned}
\end{equation*}
\begin{block}{Rate mechanism}
The finite-time feasibility bias is controlled by $r_k\downarrow0$ and the
bounded center gap.  The proof does not require a rate for
$\|\blmbd^k-\bmu^k\|$.
\end{block}
\end{frame}

\begin{frame}{Appendix: Lemma C8.3---Common Tail Index
\hfill\hyperlink{appendix:contents}{\beamerreturnbutton{Contents}}}
\scriptsize
Let $I_K:=\{\lceil K/2\rceil,\ldots,K-1\}$.  For $K\ge4$,
$|I_K|\ge K/4$.  Lemma C6.3 therefore provides one common
$\kappa\in I_K$ such that
\begin{equation*}
  \|\Delta\bx^{\kappa+1}\|^2+e_{\lambda,\kappa}^2
  +\beta_\kappa\|\bz^\kappa-\bx^{\kappa+1}\|^2
  \le\frac{4B}{K}.
\end{equation*}
Consequently,
\begin{equation*}
  \|\Delta\bx^{\kappa+1}\|+e_{\lambda,\kappa}=O(K^{-1/2}).
\end{equation*}
Because $\beta_\kappa=\Theta(K^{-1/3})$ on the tail,
\begin{equation*}
  \|\bz^\kappa-\bx^{\kappa+1}\|
  \le\sqrt{\frac{4B}{K\beta_\kappa}}=O(K^{-1/3}).
\end{equation*}
Lemma C7.2 gives
\begin{equation*}
  \|\blmbd^{\kappa+1}-\blmbd^\kappa\|=O(K^{-1/2}).
\end{equation*}
Finally, $r_\kappa=\Theta(K^{-1/3})$.  Lemmas C8.1--C8.2 now imply that
stationarity and feasibility are both $O(K^{-1/3})$ at this same index.
\end{frame}

\begin{frame}{Appendix: Main Theorem---Statement
\hfill\hyperlink{appendix:contents}{\beamerreturnbutton{Contents}}}
\hypertarget{appendix:candidate-main-theorem}{}
\small
Assume the hypotheses on $f$, $\calA$, and $h$, and choose the parameters as
stated in the Candidate algorithm.  Then for every $K\ge4$, there exists
\begin{equation*}
  \kappa\in\{\lceil K/2\rceil,\ldots,K-1\}
\end{equation*}
such that
\begin{equation*}
  \bq_\kappa^{\rm out}\in\partial h(\by_\kappa^{\rm out}),
\end{equation*}
\begin{equation*}
  \|\nabla f(\bx^{\kappa+1})
  +J_{\calA}(\bx^{\kappa+1})^\top\bq_\kappa^{\rm out}\|
  \le C_{\rm stat}K^{-1/3},
\end{equation*}
and
\begin{equation*}
  \|\calA(\bx^{\kappa+1})-\by_\kappa^{\rm out}\|
  \le C_{\rm feas}K^{-1/3}.
\end{equation*}
The constants are independent of $K$ and the requested tolerance $\eps$.
\end{frame}

\begin{frame}{Appendix: Main Theorem---Complete Proof
\hfill\hyperlink{appendix:contents}{\beamerreturnbutton{Contents}}}
\scriptsize
\noindent\textbf{Proof.}
Lemmas C1--C2 establish the exact output certificate and bounded dual orbit.
Lemmas C3--C4 give frozen three-value descent.  Lemma C5 bounds both parameter
changes by summable sequences.  Lemma C6 combines these estimates into a finite
master budget.  Lemma C8.3 selects a single tail index $\kappa$ for every term
in that budget.

At this index, Lemma C8.1 and the bounds
\begin{equation*}
  \|\Delta\bx^{\kappa+1}\|=O(K^{-1/2}),
  \qquad
  \|\bz^\kappa-\bx^{\kappa+1}\|=O(K^{-1/3})
\end{equation*}
give the stationarity estimate.  Lemmas C7.2 and C8.2, together with
\begin{equation*}
  \|\blmbd^{\kappa+1}-\blmbd^\kappa\|=O(K^{-1/2}),
  \qquad r_\kappa=O(K^{-1/3}),
\end{equation*}
give the feasibility estimate.  The subgradient relation is exact by
Lemma C8.1.  Therefore the split-KKT residual is $O(K^{-1/3})$.

Choosing $K$ so that $C K^{-1/3}\le\eps$ gives
\begin{equation*}
  K=O(\eps^{-3}).
\end{equation*}
Each iteration uses one gradient evaluation at $\bx^k$, one explicit dual
gradient, two evaluations of $\prox_{h/\rho}$, and two center averages.  Hence
the gradient and proximal-oracle complexities are both $O(\eps^{-3})$.
\hfill$\square$
\end{frame}

\end{document}
